\documentclass[11pt]{article}

\usepackage[T1]{fontenc}
\usepackage[utf8]{inputenc}
\usepackage{lmodern}
\usepackage{amsmath,amssymb,amsfonts,amsthm}
\usepackage[margin=1in]{geometry}
\usepackage{microtype}
\usepackage{xcolor}
\usepackage[colorlinks=true,linkcolor=teal,urlcolor=teal]{hyperref}

\newcommand{\F}{\mathcal F}
\newcommand{\T}{\mathcal T_1}
\newcommand{\Tr}{\operatorname{Tr}}
\newcommand{\ketbra}[1]{\lvert #1\rangle\!\langle #1\rvert}
\newcommand{\R}{\mathbb R}

\newtheorem*{lemmaD}{Lemma D.3 (rephrased)}

\title{A Short Guide to Lemma D.3\\
\large From least-noise Gaussian amplification to a conic stochastic bound}
\author{}
\date{}

\begin{document}
\maketitle

\begin{abstract}
Lemma D.3 turns a least-noise theorem for trace-preserving phase-insensitive
displacement-covariant amplifiers into an inequality valid for
trace-nonincreasing completely positive maps.  The main external input is a
stochastic-order result of Gu\c{t}\u{a} and Matsumoto.  The normalization of a
covariant subchannel, and the conversion of stochastic order into the
displayed weighted inequality, are proved in the manuscript.  This note
separates those ingredients and explains each step.
\end{abstract}

\section{What the lemma says informally}

Among phase-insensitive amplifiers with a prescribed gain, the
quantum-limited amplifier adds the least photon-number noise.  Therefore, if
one tests the output with a function that rewards low photon numbers, no
other such amplifier can score better.

Lemma D.3 also allows a probabilistic operation, whose output can have trace
less than one.  Displacement covariance forces its success probability to be
a state-independent constant $c$.  After division by $c$, the operation is
an ordinary channel to which the least-noise theorem applies; multiplication
by $c$ then gives the desired conic, or subnormalized, inequality.

\section{Background}

\subsection{One bosonic mode and thermal states}

Let $\F$ be the one-mode Fock space, with number basis
$\{\lvert k\rangle:k\geq0\}$ and number operator
\[
  \widehat N=\sum_{k\geq0}k\ketbra{k}.
\]
For $0<q<1$, the thermal state is
\begin{equation}\label{eq:thermal}
  \tau_q=(1-q)\sum_{k\geq0}q^k\ketbra{k}.
\end{equation}
Its photon number has the geometric law
$t_k=(1-q)q^k$, and its mean is $q/(1-q)$.

The Weyl displacement operators are denoted by $D(z)$, $z\in\mathbb C$.
A map $\Gamma:\T(\F)\to\T(\F)$ has amplitude gain $\sqrt\gamma$ if
\begin{equation}\label{eq:disp-cov}
  \Gamma\!\left(D(z)XD(z)^*\right)
  =D(\sqrt\gamma\,z)\Gamma(X)D(\sqrt\gamma\,z)^*.
\end{equation}
It is phase covariant if the analogous identity holds for the rotations
$e^{\mathrm i\theta\widehat N}$.  ``Phase insensitive'' is the customary
physical term for this phase symmetry.

The quantum-limited amplifier has a vacuum environment.  At intensity gain
$\gamma>1$, it maps
\begin{equation}\label{eq:thermal-output}
  \tau_q\longmapsto\tau_{q^{(\gamma)}},
  \qquad
  q^{(\gamma)}=1-\frac{1-q}{\gamma}.
\end{equation}
Indeed, its mean output occupation is
$\gamma q/(1-q)+(\gamma-1)$, which equals
$q^{(\gamma)}/(1-q^{(\gamma)})$.  Notice that
\begin{equation}\label{eq:q-order}
  q^{(\gamma)}-q
  =\frac{(\gamma-1)(1-q)}{\gamma}>0;
\end{equation}
amplification makes the thermal distribution broader.

\subsection{First-order stochastic order}

For probability laws $P,Q$ on $\mathbb Z_{\geq0}$, say that $P$ is
stochastically larger than $Q$ if
\[
  \sum_k\phi(k)P_k\geq\sum_k\phi(k)Q_k
\]
for every bounded increasing function $\phi$.  Equivalently,
\begin{equation}\label{eq:decreasing-test}
  \sum_k w(k)P_k\leq\sum_k w(k)Q_k
\end{equation}
for every bounded nonnegative decreasing function $w$.
Thus ``larger'' means that probability is shifted toward higher photon
numbers.  Useful decreasing tests include
\[
  w(k)=\boldsymbol 1_{\{k\leq K\}},
  \qquad w(k)=e^{-sk}\quad(s>0).
\]

This is an order on photon-number distributions, not an operator order
between density matrices.  Phase covariance matters because the invariant
output on the phase-invariant input $\tau_q$ is diagonal in the number basis,
so it has an unambiguous photon-number law.

\subsection{Why a covariant subchannel has constant success probability}

Suppose $\Gamma$ is completely positive and trace nonincreasing.  There is
an effect $0\leq A\leq I$ such that
\begin{equation}\label{eq:effect}
  \Tr\Gamma(X)=\Tr(AX).
\end{equation}
Equation \eqref{eq:disp-cov} and cyclicity of trace imply
$D(z)^*AD(z)=A$ for every $z$.  The Weyl representation on $\F$ is
irreducible, so its commutant consists only of scalars.  Hence
\begin{equation}\label{eq:constant-trace}
  A=cI,
  \qquad
  \Tr\Gamma(X)=c\Tr X,
  \qquad 0\leq c\leq1.
\end{equation}
This is Lemma D.1 of the manuscript.  In operational language, exact
displacement covariance prevents the probability of success from depending
on the input state.

\section{The external least-noise input}

The nontrivial fact imported from Gu\c{t}\u{a} and Matsumoto is the
following:

\medskip
\noindent
\emph{For a trace-preserving, phase-insensitive,
displacement-covariant amplifier of fixed gain, the output photon-number law
on a thermal input $\tau_q$ stochastically dominates the output law of the
vacuum-idler amplifier, namely the geometric law of
$\tau_{q^{(\gamma)}}$.}
\medskip

Their characterization represents a covariant amplifier by an idler state;
phase covariance makes the relevant idler number diagonal.  Their
least-noise lemma then identifies the vacuum idler as the stochastic
minimum.  The manuscript invokes Definition 2.2, Lemma 2.3, and the
arbitrary-gain discussion in Section III of their paper.  This classification
and least-noise theorem are \emph{not} reproved in Appendix D.

The imported statement is stronger than a comparison of mean photon number:
it controls the expectation of every bounded increasing, or equivalently
every bounded decreasing, test function.  It is exactly this stronger form
that is needed in Lemmas D.3--D.5.

\section{Statement of Lemma D.3}

\begin{lemmaD}
Let $\Gamma$ be a completely positive trace-nonincreasing one-mode map that
is displacement covariant with gain $\sqrt\gamma$ and phase covariant.  Then,
for every bounded nonnegative decreasing
$w:\mathbb Z_{\geq0}\to\R_+$,
\begin{equation}\label{eq:D3}
  \Tr\!\left[\Gamma(\tau_q)w(\widehat N)\right]
  \leq
  \Tr\Gamma(\tau_q)\,
  \Tr\!\left[\tau_{q^{(\gamma)}}w(\widehat N)\right].
\end{equation}
\end{lemmaD}

The right side is the mass of the possibly subnormalized output multiplied
by the decreasing-test expectation under the least-noise thermal output.

\section{Proof walkthrough}

\subsection*{Step 1: isolate the trace factor}

By \eqref{eq:constant-trace}, there is $c\in[0,1]$ such that
$\Tr\Gamma(X)=c\Tr X$.  In particular,
\begin{equation}\label{eq:mass-c}
  \Tr\Gamma(\tau_q)=c.
\end{equation}

If $c=0$, then $\Gamma$ is the zero map.  To see this, positivity gives
$\Gamma(X)\geq0$ and $\Tr\Gamma(X)=0$ for every $X\geq0$, hence
$\Gamma(X)=0$ on the positive cone and therefore everywhere by linearity.
Both sides of \eqref{eq:D3} are then zero.

\subsection*{Step 2: normalize the subchannel}

Assume $c>0$ and define
\[
  \Lambda=c^{-1}\Gamma.
\]
This map remains completely positive and retains both covariance properties.
Moreover,
\[
  \Tr\Lambda(X)=\Tr X,
\]
so $\Lambda$ is a trace-preserving, phase-insensitive,
displacement-covariant amplifier of gain $\sqrt\gamma$.

\subsection*{Step 3: apply the imported stochastic-order theorem}

Let $P$ be the photon-number law of $\Lambda(\tau_q)$ and let
\[
  Q_k=(1-q^{(\gamma)})(q^{(\gamma)})^k
\]
be the law of $\tau_{q^{(\gamma)}}$.  The Gu\c{t}\u{a}--Matsumoto result
says that $P$ is stochastically larger than $Q$.  Applying the equivalent
decreasing-test formulation \eqref{eq:decreasing-test} gives
\begin{equation}\label{eq:normalized-bound}
  \Tr\!\left[\Lambda(\tau_q)w(\widehat N)\right]
  \leq
  \Tr\!\left[\tau_{q^{(\gamma)}}w(\widehat N)\right].
\end{equation}

\subsection*{Step 4: restore the trace factor}

Multiplying \eqref{eq:normalized-bound} by $c$ and using
$\Gamma=c\Lambda$ and \eqref{eq:mass-c} yields precisely \eqref{eq:D3}.
No additional Gaussian calculation is needed.

\section{How to read the inequality}

Taking $w=1$ gives equality $c=c$.  Taking
$w(k)=\boldsymbol 1_{\{k\leq K\}}$ gives
\[
  \Pr_{\Gamma,\mathrm{sub}}\{N\leq K\}
  \leq c\,\Pr_{\tau_{q^{(\gamma)}}}\{N\leq K\}.
\]
After dividing by $c$ when $c>0$, the left side becomes the low-energy tail
probability of $\Lambda(\tau_q)$.  The inequality says that even a
probabilistic covariant operation cannot conditionally produce a quieter
amplified output than the vacuum-idler amplifier.

The specific weight later needed for fidelity is
\[
  w_k=\sqrt{\frac{(1-q)q^k}
                       {(1-q^{(\gamma)})(q^{(\gamma)})^k}}
  =\sqrt{\frac{1-q}{1-q^{(\gamma)}}}
   \left(\frac{q}{q^{(\gamma)}}\right)^{k/2}.
\]
By \eqref{eq:q-order}, this is bounded, nonnegative, and decreasing, so it
falls exactly within Lemma D.3.

\section{What is imported and what is proved}

Lemma D.1 internally proves $\Tr\Gamma(X)=c\Tr X$ from Weyl
irreducibility.  Lemma D.3 proves the normalization by $c$ and converts the
decreasing-test definition into \eqref{eq:D3}.  The sole imported input is
Gu\c{t}\u{a}--Matsumoto's stochastic minimality of the vacuum-idler output.

\section{Dependencies and role in the paper}

The logical chain is
\[
  \text{Lemma D.1}
  +\text{Gu\c{t}\u{a}--Matsumoto least-noise theorem}
  \longrightarrow \text{Lemma D.3}.
\]
Lemma D.3 is then applied one mode at a time in Lemma D.4.  Lemma D.4
controls the tangent weight used in Lemma D.5, which gives the multimode
fidelity upper bound in Theorem 3.4.  The same tensorized inequality is also
applied branchwise to the residual subinstrument in the hybrid Gaussian
product bound.

The hypothesis of phase covariance is normally arranged by phase twirling.
At the origin, both the thermal input and thermal target are phase invariant,
and concavity of root fidelity ensures that this twirling does not decrease
the payoff.

\section{Minimal prerequisite checklist}

To follow the internal proof, it is enough to know completely positive maps,
Weyl covariance, irreducibility and commutants, and first-order stochastic
dominance.  The specialized ingredient is the Gu\c{t}\u{a}--Matsumoto
least-noise theorem, which Appendix D uses as a black box.

\medskip
\noindent
\textbf{Primary source.}
M. Gu\c{t}\u{a} and K. Matsumoto,
``Optimal cloning of mixed Gaussian states,''
\emph{Physical Review A} \textbf{74}, 032305 (2006),
\href{https://arxiv.org/abs/quant-ph/0605161}{arXiv:quant-ph/0605161}.

\end{document}
